\subsection{Preliminary results}

\paragraph{$n$-friends.}
We first tested the generalized search on the knot $6_2$.  Searching all
geodesics of length at most $3$ in the relevant filled manifold found one
$1$-friend.  SnapPy identifies its exterior as \texttt{K13n3596}, in
agreement with the example used to test the implementation.  The analogous
$3$-surgery search found the same exterior from two different geodesic words;
SnapPy identifies it as \texttt{K14n10164}.  In the orientation used by the
census, its $-3$-surgery agrees with the $3$-surgery of $6_2$.  We therefore
mirror the recovered diagram, obtaining a knot whose positive $3$-surgery
agrees with that of $6_2$.  The resulting $1$- and $3$-friend diagrams have
$13$ and $14$ crossings, respectively, and are saved as
\texttt{6\_2\_1\_friend\_K13n3596.lnk} and
\texttt{6\_2\_3\_friend\_K14n10164\_mirror.lnk}.

We also searched for $1$-friends of the Conway knot, represented in SnapPy by
\texttt{K11n34}.  With the same length cutoff, the program returned five
geodesic hits, with exterior volumes approximately $12.90272$, $13.70618$,
$13.81844$, and $14.30116$ (the last volume occurred twice).  These exteriors
are not identified with named knots in the installed SnapPy census.  We
recovered and saved diagrams for the two lowest-volume examples.  After
choosing the mirror when necessary so that positive $1$-surgery agrees with
positive $1$-surgery on Conway, these diagrams have $34$ and $40$ crossings;
they are denoted \texttt{conway\_1\_friend\_A.lnk} and
\texttt{conway\_1\_friend\_B.lnk}.  For each example we also save a
\texttt{.tri} file, which records the exterior independently of the chosen
diagram.  The four recovered knot diagrams are shown in
Figure~\ref{fig:preliminary-n-friends}.

\begin{figure}[htbp]
  \centering
  \begin{minipage}[t]{0.47\textwidth}
    \centering
    \includegraphics[width=\linewidth]{images/6_2_1_friend_K13n3596.png}
    \small $1$-friend of $6_2$: \texttt{K13n3596}
  \end{minipage}\hfill
  \begin{minipage}[t]{0.47\textwidth}
    \centering
    \includegraphics[width=\linewidth]{images/6_2_3_friend_K14n10164_mirror.png}
    \small $3$-friend of $6_2$: the mirror of \texttt{K14n10164}
  \end{minipage}

  \vspace{0.75em}

  \begin{minipage}[t]{0.47\textwidth}
    \centering
    \includegraphics[width=\linewidth]{images/conway_1_friend_A.png}
    \small Conway $1$-friend A
  \end{minipage}\hfill
  \begin{minipage}[t]{0.47\textwidth}
    \centering
    \includegraphics[width=\linewidth]{images/conway_1_friend_B.png}
    \small Conway $1$-friend B
  \end{minipage}
  \caption{Knot diagrams recovered by the $n$-friend search.  Breaks at
  crossings indicate the underpassing strand.}
  \label{fig:preliminary-n-friends}
\end{figure}

\paragraph{$n$-RBG links.}
We next used these pairs as input to the generalized RBG search.  For the pair
$6_2$ and \texttt{K13n3596}, the program recovered an $11$-crossing
$1$-special RBG link.  This is the small example from Qin et al.: with the
components ordered red, blue, green, its surgery slopes are
$(-1,1),(0,1),(0,1)$.  Both the red--blue and red--green sublinks are Hopf
links, every component is unknotted, and the framed linking matrix has
determinant $-1$.  Filling the red and green components recovers the exterior
of $6_2$, while filling the red and blue components recovers the exterior of
\texttt{K13n3596}.  The recovered diagram is saved as
\texttt{6\_2\_K13n3596\_1\_RBG.lnk}.

The same experiment also produced a small RBG presentation for Conway and
\texttt{conway\_1\_friend\_A}.  Searching blue geodesics to radius $3$ and
dual curves with at most $16$ segments returned an $18$-crossing $1$-special
RBG diagram with the same three framings.  Direct filling checks recover the
Conway exterior and the saved friend exterior, and the Hopf-link and
determinant tests again pass.  This diagram is saved as
\texttt{conway\_friend\_A\_1\_RBG.lnk}.  Thus the initial experiments supply
both individual $n$-friend diagrams and explicit RBG diagrams realizing two
of the common-surgery pairs.  Figure~\ref{fig:preliminary-rbg-links} shows
the two RBG diagrams, with their components colored red, blue, and green.

\begin{figure}[htbp]
  \centering
  \begin{minipage}[t]{0.47\textwidth}
    \centering
    \includegraphics[width=\linewidth]{images/6_2_K13n3596_1_RBG.png}
    \small The $1$-RBG link for $6_2$ and \texttt{K13n3596}
  \end{minipage}\hfill
  \begin{minipage}[t]{0.47\textwidth}
    \centering
    \includegraphics[width=\linewidth]{images/conway_friend_A_1_RBG.png}
    \small The $1$-RBG link for Conway and friend A
  \end{minipage}
  \caption{The two $1$-special RBG links found in the preliminary search.
  Components are displayed in red--blue--green order.}
  \label{fig:preliminary-rbg-links}
\end{figure}
